\chapter{RELATED WORK AND BACKGROUND}

\section{Shortest-Path Algorithms for Conversion Path Optimization}

Dijkstra~\citep{dijkstra1959} introduced one of the most fundamental shortest-path algorithms for weighted graphs with non-negative edge weights. The algorithm incrementally explores nodes in order of increasing accumulated cost using a priority queue, guaranteeing optimality when all edge weights are non-negative. When implemented with a binary min-heap, it achieves a time complexity of $O(|E|\log|V|)$~\citep{cormen2009}. A theoretically superior bound of $O(|E| + |V|\log|V|)$ is achievable using Fibonacci heaps~\citep{fredman1987fibonacci}, though the practical overhead of Fibonacci heaps often makes binary-heap implementations preferable for moderately sized graphs. In the context of this project, Dijkstra's algorithm serves as the baseline, applied to a log-transformed probability graph where edge weights represent the negative log-probability of transitions. Its well-understood complexity and correctness guarantees make it the natural choice for baseline comparison and scalability analysis.

\section{Heuristic and Approximate Shortest-Path Approaches}

While Dijkstra's algorithm guarantees optimality, it explores the full reachable graph from the source. Informed search algorithms such as A*~\citep{hart1968astar} reduce exploration by using an admissible heuristic function to guide node expansion toward the target, achieving sub-Dijkstra runtime when a valid heuristic exists. Pearl~\citep{pearl1984heuristics} provides a comprehensive theoretical treatment of heuristic search, establishing formal conditions under which informed search algorithms guarantee optimal solutions. In probabilistic graph settings, however, designing a tight and admissible heuristic over log-probability weights is non-trivial, making heuristic approaches difficult to apply directly. Beam search offers an alternative approximation strategy by limiting the frontier to the top-$k$ most promising candidates at each step~\citep{cormen2009}. Although beam search can be very fast in practice, it provides no optimality guarantee and may miss the globally optimal path entirely. This project adopts a different approximation strategy---probability-threshold pruning---which is deterministic, parameter-controlled, and guarantees convergence to the optimal solution as the threshold approaches zero, offering a more principled trade-off than beam search. Unlike A*, which requires a problem-specific admissible heuristic that is difficult to define over log-probability edge weights, and unlike beam search, which discards nodes based on rank with no convergence guarantee, the proposed pruning operates directly on accumulated path cost against a single interpretable parameter $\tau$, requiring no domain-specific heuristic design and providing a formal convergence proof as $\tau \to 0$.

Another important family of shortest-path speedup methods is \emph{Contraction Hierarchies}, which accelerates queries by preprocessing the graph and adding shortcut edges to preserve shortest-path distances. Such methods are highly effective for static road-network-style graphs, but they are less naturally suited to the probabilistic customer-journey setting considered here, where transition weights are derived from behavioral data and experimental emphasis is placed on controllable approximation rather than heavy preprocessing. Nevertheless, Contraction Hierarchies provide an important point of comparison within the broader landscape of shortest-path acceleration methods.

\section{Markov Chain Models for Probabilistic Navigation}

The probabilistic transition model used in this project---where each user state transitions to the next with a fixed probability---is formally equivalent to a first-order Markov chain~\citep{norris1998markov}. Markov chains have been widely used to model sequential user behaviour in web navigation, recommendation systems, and clickstream analysis~\citep{rabiner1989hmm,baum1966statistical}. The maximum-probability path problem on a Markov chain is structurally identical to the Viterbi algorithm~\citep{viterbi1967error}, which finds the most likely state sequence in a Hidden Markov Model using dynamic programming in $O(|V|^2 T)$ time for $T$ time steps. However, the Viterbi algorithm is designed for sequences of fixed length, whereas the present problem involves finding the most probable \emph{path} from $s$ to $t$ in a general directed graph without a fixed path length. The log-probability shortest-path formulation adopted here generalizes the Viterbi principle to arbitrary directed graphs and applies it efficiently using Dijkstra's algorithm, which has better practical complexity on sparse graphs.

\section{Algorithmic Foundations of Efficient Graph Traversal}

Cormen et al.~\citep{cormen2009} provided a comprehensive treatment of graph algorithms, including shortest-path search, graph traversal, and priority-queue-based optimization techniques. The textbook formally analyzes time and space complexity, establishes correctness proofs for Dijkstra's algorithm, and discusses data structure choices including binary heaps and adjacency lists. These foundations directly support the design decisions in this project: the use of adjacency lists for $O(|V|+|E|)$ space, binary heaps for $O(\log|V|)$ extract-min and insert operations, and hash maps for $O(1)$ distance lookup. The log-probability weight transformation exploits the monotone relationship between probability maximization and shortest-path minimization, a standard technique in probabilistic graph optimization~\citep{manning2008nlp}.

\section{Sequential Pattern Mining for Frequent Journey Discovery}

Pei et al.~\citep{pei2001prefixspan} proposed PrefixSpan, a pattern-growth-based algorithm for mining frequent sequential patterns without candidate generation. By projecting the database based on frequent prefixes, PrefixSpan significantly reduces search space compared to Apriori-based approaches and demonstrated strong scalability on large sequence datasets. In the context of this project, PrefixSpan is conceptually relevant as a method for identifying frequently taken customer journey paths without enumerating all possible sequences. Its prefix-based pruning principle shares conceptual similarities with the threshold-based pruning strategy proposed in this work, where low-probability path prefixes are discarded early to reduce the effective search space.

\section{Survey of Frequent and Sequential Pattern Mining Techniques}

Han et al.~\citep{han2007sequential} presented a comprehensive survey of frequent pattern mining methods, covering association rules, sequential patterns, and constrained pattern mining. The survey highlights key challenges including scalability, pattern explosion, and the need for efficient pruning strategies. These challenges closely mirror those encountered in customer journey analysis, where the number of simple paths from source to target grows combinatorially with graph size. This project draws on the general principle of threshold-based pruning to control computational complexity, filtering out low-probability paths before full expansion in a manner analogous to minimum support thresholds in frequent pattern mining.

\section{Modeling User Navigation as Directed Graphs}

Bucklin and Sismeiro~\citep{bucklin2003clickstream} modeled user browsing behaviour using clickstream data, representing navigation patterns as structured sequences and probabilistic transition graphs. Their empirical analysis demonstrated that user navigation can be effectively captured using probabilistic transition models, providing insights into user behaviour and engagement at scale. Liu~\citep{liu2011webmining} further surveyed web usage mining techniques---including clickstream analysis, hyperlink analysis, and content mining---providing a broader context for how navigation data is processed into actionable graph representations. This work directly supports the modeling approach adopted in this project, where customer journeys are represented as directed weighted graphs derived from interaction sequences, and transition probabilities are estimated from session frequency counts using maximum likelihood estimation.

\section{Probabilistic Models for Navigation and Path Selection}

Resnick and Varian~\citep{resnick1997recommender} discussed probabilistic models of user behaviour, including transition-based navigation models that capture the likelihoods of moving between system states. Such models are commonly used to estimate user preferences and navigation patterns in large-scale recommendation and information retrieval systems. The probabilistic transition modeling approach adopted in this project---assigning edge weights as negative log-probabilities and finding the minimum-weight path---is consistent with the broader class of log-linear models used in language modeling and information retrieval~\citep{manning2008nlp,manning2008ir}, where the negative log-likelihood serves as a natural additive cost function.

\section{Scalable Graph Processing and Large-Scale Navigation Analysis}

Recent advances in large-scale graph processing systems further emphasize the importance of efficient algorithmic design when analyzing massive navigation graphs. Malewicz et al.~\citep{malewicz2010pregel} introduced Pregel, a vertex-centric computational model for scalable graph processing, demonstrating how distributed systems can efficiently handle graphs containing millions or billions of edges. Subsequent research in large-scale graph analytics has examined optimization strategies for traversal, pruning, and distributed processing to reduce computational overhead in dense and evolving networks~\citep{sahu2020scaling}. These studies reinforce the practical relevance of designing efficient path-search algorithms and pruning strategies when working with large customer journey graphs, where computational complexity grows rapidly with graph size. The present project builds upon these principles by evaluating algorithmic efficiency and search-space reduction techniques within a single-machine, controlled experimental framework, where the goal is to characterize complexity empirically rather than to achieve distributed scalability.

\vspace{0.5cm}

\noindent The reviewed literature establishes a strong theoretical and methodological foundation for this project. Dijkstra's algorithm~\citep{dijkstra1959} and its complexity analysis~\citep{cormen2009,fredman1987fibonacci} provide the basis for the baseline implementation and complexity comparison. The Markov chain and Viterbi frameworks~\citep{norris1998markov,viterbi1967error,rabiner1989hmm} situate the probabilistic path problem in a well-studied theoretical context. Heuristic search methods~\citep{hart1968astar,pearl1984heuristics} and approximate approaches establish the landscape of alternatives against which the pruning strategy is positioned. Sequential pattern mining techniques~\citep{pei2001prefixspan,han2007sequential} and clickstream modeling~\citep{bucklin2003clickstream,liu2011webmining} validate the representation and motivate pruning as a scalability mechanism. Probabilistic modeling literature~\citep{resnick1997recommender,manning2008nlp} supports the log-probability transformation. Scalable graph systems~\citep{malewicz2010pregel,sahu2020scaling} contextualize the scalability challenges addressed. The Erd\H{o}s--R\'{e}nyi model~\citep{erdosrenyi1960} grounds the synthetic data generation methodology, and real-world benchmarks~\citep{retailrocket2015,uci2008eshop} validate experimental generalizability.

\section{Graph Cycles in Customer Journey Navigation}

A critical modeling question is whether the customer journey graph contains cycles---that is, whether a user may return to a previously visited state, such as \textit{Landing $\to$ Browse $\to$ Product $\to$ Browse}. Real-world clickstream data frequently exhibits such revisitation behaviour~\citep{bucklin2003clickstream}: users loop between states (e.g., browsing multiple product pages before adding to cart) before eventually converting or abandoning the session. The customer journey graph is therefore modeled as a general directed weighted graph that \emph{may contain cycles}.

A key concern is whether Dijkstra's algorithm remains correct on cyclic graphs. Cormen et al.~\citep{cormen2009} establish that Dijkstra's algorithm is correct on any directed graph with non-negative edge weights, regardless of whether cycles are present. The log-probability transformation $w(u,v) = -\log\,p(u,v) \geq 0$ guarantees non-negativity for all $p(u,v) \in (0,1]$, so Dijkstra's algorithm is directly applicable to cyclic navigation graphs without any cycle-detection preprocessing. Since each cycle traversal adds strictly positive cost, the algorithm naturally avoids cycling: once a node is finalized, it is never re-extracted from the priority queue. This property means no modifications to the standard algorithm are required to handle the cyclic structure that is inherent in real user navigation data.

\section{Synthetic Graph Generation for Controlled Evaluation}

To evaluate algorithmic performance in a controlled and reproducible manner, synthetic graphs are used as the primary experimental dataset. The generation methodology is grounded in the Erd\H{o}s--R\'{e}nyi random graph model~\citep{erdosrenyi1960}, which provides a theoretically principled basis for generating random directed graphs with predictable structural properties (connectivity, degree distribution) at a given size and density.

Graphs are generated as follows. Each node represents an interaction state. For each graph instance with $|V|$ nodes and target average out-degree $\bar{d} \in \{2, 5, 10\}$, directed edges are sampled independently such that the expected number of edges is $\bar{d} \cdot |V|$. Transition probabilities are assigned under two distributions: (i) \textbf{uniform} $\mathcal{U}(0.1, 1.0]$, simulating balanced navigation, and (ii) \textbf{skewed power-law} with exponent $\alpha = 2$, where a small number of transitions carry high probability and the majority carry low probability, reflecting the heterogeneous click patterns observed in real e-commerce data~\citep{bucklin2003clickstream,retailrocket2015}. For each node, outgoing probabilities are normalized to $\sum_v p(u,v) \leq 1$, ensuring a valid sub-stochastic transition system. A fixed source $s$ and target $t$ are assigned per instance, with at least one valid $s$-to-$t$ path guaranteed by a depth-first seeding step. All experiments use fixed random seeds for full reproducibility.

Synthetic generation is preferred as the primary substrate because it allows independent control of $|V|$, $|E|$, and the probability distribution---three variables that cannot be decoupled in observational data. Real-world datasets (RetailRocket~\citep{retailrocket2015}, UCI E-Commerce Clickstream~\citep{uci2008eshop}) are used as secondary validation to confirm findings generalize beyond synthetic graphs.

\vspace{0.5cm}
\noindent\textbf{Research Gap.} Despite the extensive literature above, existing studies largely treat shortest-path computation~\citep{dijkstra1959,cormen2009}, probabilistic modeling~\citep{norris1998markov}, and sequential pattern mining~\citep{pei2001prefixspan} as separate problems. Limited work integrates all three within a unified framework tailored to customer conversion path analysis, with explicit treatment of cyclic navigation graphs, principled synthetic data generation~\citep{erdosrenyi1960}, and empirical characterization of the time--optimality trade-off of threshold-based pruning.

The present project addresses this gap by combining probabilistic transition modeling, threshold-controlled pruning, and Dijkstra's algorithm into a cohesive experimental framework. By systematically evaluating performance across diverse synthetic and real-world graph configurations---including cyclic graphs---the study bridges classical graph algorithms, Markov chain theory, and practical navigation graph analysis within a single experimental design. This work contributes a threshold-based pruning strategy integrated with Dijkstra's algorithm for probabilistic navigation graphs, evaluated through empirical complexity analysis and an explicit time--optimality trade-off characterization.
